\documentclass{situationist-preprint}

\title{The Disentangling Argument Reconsidered}
\msid{MTE-15-b}
\msfamily{MTE-15}
\msdiscipline{Metaethics}
\mstype{Companion manuscript}
\mskeywords{thick concepts; disentangling argument; shapelessness; McDowell; separabilism; outsider}

\begin{document}
\maketitle

\begin{abstract}
The disentangling argument, due to McDowell, holds that thick concepts cannot be split into a descriptive and an evaluative component, because an outsider who does not share the evaluative perspective could not follow the concept's extension. Critics reply that the argument trades on an ambiguity between sharing a perspective and understanding it, and that shapelessness relative to a vocabulary is not shapelessness in nature. We reconsider the argument in the form MTE-15-a gives it and show that both replies fail: the outsider of the argument is one who neither shares nor understands the evaluation, and the shape that matters is the shape that can be mastered, not the shape that exists.
\end{abstract}

\section{The argument and two replies}

\citet{McDowell1981} argued that the extension of a thick concept need not correspond to any kind recognizable without the evaluative interest that informs it, so that a person who stood outside that interest could not learn to go on applying the concept to new cases. If this is right, the concept cannot be disentangled into a descriptive part that fixes the extension and an attitude that is attached to it. MTE-15-a states the argument with a definition of descriptive concepts as those masterable without sharing or understanding any evaluative attitude (Definition~2.1 of MTE-15-a), a factoring thesis that includes a condition on competence (Definition~2.2(b)), and the premise of shapelessness (Principle~P2), and proves that \emph{lewd} does not factor (Theorem~3.2).

Two replies are prominent. The first is that the argument shows at most that one must \emph{understand} the evaluative point of the concept, which an anthropologist can do without endorsing it; \citet{Williams1985} describes such an observer, and \citet{Blackburn1992} takes the point to leave separabilism intact. The second is that the extension of a thick concept may be shapeless relative to our descriptive vocabulary and yet coincide with some descriptive property, perhaps a complex disjunction, which the separabilist can take as the descriptive component.

\section{The replies examined}

\paragraph{Understanding without sharing.} The first reply concedes the conclusion. Factoring requires that mastery of the thick concept consist in mastery of a descriptive concept together with a disposition to evaluate (Definition~2.2(b) of MTE-15-a). An observer who can project \enquote{lewd} only by grasping what its users find offensive and why has not mastered a descriptive concept, since by Definition~2.1 a descriptive concept is one that can be mastered without understanding any evaluative attitude. Her competence includes an understanding of the evaluation, and so it does not decompose as factoring requires. The insightful observer is therefore not a counterexample to Theorem~3.2 of MTE-15-a but an illustration of it: what she needs in order to go on is exactly the evaluative component, held in imagination rather than in endorsement. The distinction between sharing and understanding matters for the question whether one can use a thick concept without being committed to its evaluation, but not for the question whether the concept factors.

\paragraph{Shape in nature.} The second reply misidentifies the premise. Principle~P2 of MTE-15-a is not that there is no descriptive property coextensive with \emph{lewd}; it is that there is no descriptive \emph{concept} with that extension across the cases competent users would classify. A descriptive concept is individuated by what mastery of it requires. A disjunction of descriptive conditions, one for each case that users of \enquote{lewd} would classify, is a property but not a concept that can be learned by someone without the evaluation, since nothing in the disjuncts tells the learner how to extend the list to the next case. Supervenience of the thick on the descriptive, which the inseparabilist accepts, guarantees that such a property exists; it does not provide the separabilist with a descriptive component in the sense that factoring requires.

\section{What the argument does not show}

The disentangling argument, so reconsidered, does not show that thick concepts are not world-guided, nor that the facts they report are not natural facts; the shapeless property is a natural property, and those who share the evaluation are guided by the world in applying it. Nor does it show that disagreements over the application of a thick term are disagreements in attitude only. A community that shares the evaluative point of \enquote{lewd} and applies it differently misapplies the same concept, as MTE-13 holds for moral terms generally. Expressivism must, by MTE-15-a, give up the factoring of thick concepts, but this leaves open the expressivist semantics of MTE-03, which does not depend on factoring.

\section{Conclusion}

The disentangling argument survives its two main replies. The outsider it requires is one who neither shares nor understands the evaluation, and an observer who understands without sharing confirms the argument rather than refuting it. The shapelessness it requires is a fact about the concepts that can be mastered, not about the properties that exist. Thick concepts do not factor into description and evaluation.

\begin{thebibliography}{9}\small
\bibitem[Blackburn(1992)]{Blackburn1992} Blackburn, S. (1992). Morality and thick concepts: Through thick and thin. \emph{Proceedings of the Aristotelian Society, Supplementary Volume}, 66, 285--299.
\bibitem[McDowell(1981)]{McDowell1981} McDowell, J. (1981). Non-cognitivism and rule-following. In S. Holtzman and C. Leich (eds.), \emph{Wittgenstein: To Follow a Rule}, 141--162. Routledge and Kegan Paul.
\bibitem[Williams(1985)]{Williams1985} Williams, B. (1985). \emph{Ethics and the Limits of Philosophy}. Fontana.
\end{thebibliography}

\end{document}
